\begin{pastexam}[label={px:ncku-111-ee-ds-04}]{成大電機所 111 學年度 單選題第 4 題}
（單選；每題 4 分）

A \underline{\hspace{2em}} is a data structure that organizes data similar to a line in the supermarket, where
the first one in line is the first one out.

A.~queue linked list \qquad B.~stacks linked list \qquad C.~both of them \qquad D.~neither of them
\end{pastexam}

\begin{pastsoln}
答案是 A。

超市排隊，先排的先結帳，就是\textbf{先進先出}（FIFO），也就是佇列（資結書第~6.2~節）。 選項的「queue linked list」
是用串列實作的佇列：從尾端加入、從前端移除，兩邊各記一個指標，兩個操作都是 $\Theta(1)$。
\begin{itemize}
\item A.~佇列：對。
\item B.~堆疊：後進先出，最後排的先出來，不像排隊。 錯。
\item C.~兩者都是：堆疊不是。 錯。
\item D.~兩者都不是：佇列就是。 錯。
\end{itemize}
\end{pastsoln}
